% Standalone section for problems/3.2/proof.tex — \input this file.
% Written 2026-08-01 (life window, Fable campaign, phase apparition).
% Companion: FABLE_SECTION_value_distribution.tex. Audit scripts named inline
% (all in research/scripts/). Full session log: FABLE_NOTES_energy_bootstrap.md.

\section{The mod-$p$ structure theory: strong reflection, the Frobenius tower,
and the apparition theorem}
\label{sec:apparition-tower}

Throughout $p\ge7$, $b_n=\sum_k\binom nk^2\binom{n+k}k^2$,
$F(t)=\sum_{n\ge0}b_nt^n$, $q(t)=1-34t+t^2$,
$A_p(t)=\sum_{n<p}b_nt^n$, $\chi_p=\bigl(\tfrac{-6}p\bigr)$,
$P(m)=34m^3+51m^2+27m+5$.

\subsection{Strong reflection}

\begin{theorem}[strong reflection]\label{thm:strong-reflection}
Every solution germ $y$ of the Ap\'ery recurrence mod $p$ on the regular
window satisfies $y_{p-1-n}=y_n$.  In particular $A_p$ is palindromic,
$b_{p-1}\equiv1$, $b_{p-2}\equiv5$.
\end{theorem}

\begin{proof}
At $c\equiv-\tfrac12$: $P(c)=0$ (the factor $2m+1$) and $(c+1)^3=-c^3$, so
the recurrence at $m=c$ degenerates to $y_{c+1}=y_{c-1}$ for every solution.
The reflected sequence $\tilde y_n=y_{p-1-n}$ is a solution
($P(-1-n)=-P(n)$) agreeing with $y$ at the consecutive indices $c,c+1$;
propagate.  (Machine check: 80/80 random germs, \texttt{strong\_reflection.py}.)
\end{proof}

Corollaries used below: the mirror root $\Psi_{h,k}\bigl(\tfrac{p-1-h-k}2\bigr)=0$
for $h\equiv k\ (2)$, and the universal even-lag type-II centers
$A_d\bigl(-\tfrac{d+1}2\bigr)=1$, $B_d\bigl(-\tfrac{d+1}2\bigr)=0$
(\S\ref{sec:value-distribution}; \texttt{mirror\_identity.py},
\texttt{factorial\_ledger.py}).

\subsection{The complete mod-$p$ structure theorem}

\begin{theorem}[structure]\label{thm:structure}
Over $\mathbb F_p$ the truncated generating polynomial factors as
\[
A_p=\begin{cases} S_p^2,&\chi_p=+1,\\ q\cdot S_p^2,&\chi_p=-1,\end{cases}
\qquad S_p(0)=1,
\]
and $S_p$ is the truncation of a FIXED $\mathbb Q$-series:
$S_p=[\sqrt F]_{(p-1)/2}$ when $\chi_p=+1$, and
$S_p=[\sqrt{F/q}]_{(p-3)/2}$ when $\chi_p=-1$.  The fixed branches
$\tau=\sqrt F$, $\sigma=\sqrt{F/q}$ satisfy second-order half-integer-index
recurrences (e.g.\ $4(j+2)^2\tau_{j+2}=2(68j^2+170j+107)\tau_{j+1}-(2j+1)^2\tau_j$)
and are algebraic pullbacks of a single classical ${}_2F_1$ with exact
coefficient formulae by Lagrange inversion
(\texttt{CODEX\_JACOBSTHAL\_DEEP.md}).  Consequently the entire mod-$p$
structure of the Ap\'ery sequence is carried by: three universal sequences
($D_r$, $E_r$, $\tilde U_r$ of the tower below), two fixed $\mathbb Q$-series
($\tau$, $\sigma$), and exactly two scalars per prime ($\chi_p$ and $\beta_p$).
(Square-factor identification verified coefficientwise $p<150$ and
spot-checked in both classes: \texttt{sp\_spotcheck.py}.)
\end{theorem}

\subsection{The Frobenius tower and the master corner identity}

\begin{theorem}[tower]\label{thm:tower}
For $0\le r<p$ the two-digit block expansion
\[
b_{p+r}\equiv 5b_r+10p\,D_r+p^2E_r+p^3(\tilde U_r+\beta_p b_r)
+p^4(\hat U_r-6W^a_r+2\beta_pD_r+8a_r+y_4b_r)\pmod{p^5}
\]
holds with UNIVERSAL (p-independent) sequences at orders $p^0$–$p^2$, where
$D_r=\sum_k\binom rk^2\binom{r+k}k^2(H_{r+k}-H_{r-k})$.  (Attribution: the
$p^0$ layer is Gessel's congruence (1982); the $p^1$--$p^2$ layers are the
$n=1$ case of A.~Straub, arXiv:2301.12248, Monatsh.\ Math.\ (2024),
Thm.~1.3.  New here: the $p^3$/$p^4$ universality, the closed forms of the
invariants, and the all-order tower structure.)  The only
$p$-dependent invariants through order $p^4$ are $\beta_p=(b_p-5)/p^3\bmod p$
and $y_4$ (second invariant, weight-3 Kummer quotient).  Moreover:
\[
\textbf{(master corner identity)}\qquad
b_p-5\;\equiv\;-7\,p^2\,H^{(2)}_{p-1}\pmod{p^5},
\]
which subsumes Beukers' $b_p\equiv5\ (p^3)$, the identity
$\beta_p\equiv-\tfrac{14}3B_{p-3}\ (p)$ (theorem status: direct
binomial/harmonic proof, and Thm~1.1 of Ji-Cai Liu, arXiv:2404.16636,
specializes to it), and the $y_4$ Kummer-quotient formula.  In particular
$\beta_p=0\iff p$ is a Wolstenholme prime, and on the divisibility locus
$p\mid b_r$ the terms $\beta_pb_r$, $y_4b_r$ self-cancel: the tower is
universal through $p^4$ at bad primes.
Unconditional proof (cron, Q6330): the exact block identity
$b_{N+r}=b_N F_r(N)-N^3b_{N-1}G_r(N)$ (recurrence uniqueness) plus the two
endpoint lemmas $b_p\equiv5\ (p^3)$, $b_{p-1}\equiv1\ (p^2)$ assemble the
tower through $p^4$; the $N^3$ factor is the entire rigidity mechanism, and
all $p$-dependence concentrates in the two endpoint digit sequences
(all-order form).  (Verification: supercongruence 63/63
\texttt{cron\_supercong\_check.py}; $\beta_p$ 31/31
\texttt{beta\_bernoulli\_check.py} — equivalently $\beta_p\equiv
7\cdot(\text{Wolstenholme quotient})$; master identity 21/21
\texttt{harmonic\_p5\_check.py}; independent ledgers 25--27/27.)
\end{theorem}

\subsection{The apparition theorem}

\begin{theorem}[Ap\'ery apparition law]\label{thm:apparition}
Let $e=\tfrac{p-1}2$ and let $S_p$ be the branch polynomial of
Theorem~\ref{thm:structure}, $D=\deg S_p\in\{e,e-1\}$.
\begin{itemize}
\item[(i)] \emph{(endpoint)} $\operatorname{lc}(S_p)\equiv\bigl(\tfrac{-2}p\bigr)$;
\item[(ii)] \emph{(branch reflection)}
$t^{D}S_p(1/t)=\bigl(\tfrac{-2}p\bigr)S_p(t)$, i.e.\
$s_{D-j}\equiv\bigl(\tfrac{-2}p\bigr)s_j$;
\item[(iii)] \emph{(quarter-point law)} the central coefficient vanishes iff
the center is integral and $\bigl(\tfrac{-2}p\bigr)=-1$:
\[
\tau_{(p-1)/4}\equiv0 \iff p\equiv5\ (24),\qquad
\sigma_{(p-3)/4}\equiv0 \iff p\equiv23\ (24),
\]
and in no other residue class mod 24 does the floor-quarter coefficient of
the relevant branch vanish.
\end{itemize}
\end{theorem}

\begin{proof}[Proof (two independent routes)]
\emph{Route A (Franel pullback endpoint).}  With $f_n=\sum_k\binom nk^3$,
$h(u)=\sum f_nu^n$, $\phi(u)=u(1-8u)/(1+u)$, the exact identity
$F(\phi(u))=(1+u)h(u)^2$ (characteristic zero; CFVZ) together with the Dwork
congruence $F=A_pF(t^p)$ and the Lucas factorization $h=H_ph(u^p)$ gives
$A_p(\phi(u))=H_p(u)^2/(1+u)^{p-1}$ with $H_p$ monic of degree $p-1$
($f_{p-1}\equiv1$); in the $\chi_p=-1$ branch additionally
$q(\phi(u))=\bigl((1-16u-8u^2)/(1+u)\bigr)^2$.  Identifying $S_p$ with the
fixed-branch truncation (unit argument in $\mathbb F_p[[t]]/(t^p)$) and
comparing leading terms at $u=\infty$ yields
$\operatorname{lc}(S_p)=(-8)^{-e}=\bigl(\tfrac{-2}p\bigr)$ in both branches;
palindromy of $A_p$ (Theorem~\ref{thm:strong-reflection}) and unique
factorization give (ii), and (ii) forces (iii) at the center.
(Every ingredient machine-verified: \texttt{r18\_proof\_check.py};
endpoint 165/165 \texttt{leading\_coeff\_check.py}; reflection 92/92
\texttt{branch\_reflection\_check.py}; quarter table $p<2000$
\texttt{full24\_jacobsthal.py}.)

\emph{Route B (direct reversal).}  An independent proof of (ii) without the
endpoint evaluation, via the reversal symmetry of the branch polynomial, is
given in \texttt{CODEX\_JACOBSTHAL\_DEEP.md} (commits d9af2b9--77e089b), with
its own verification ledger (all $p<3000$).
\end{proof}

\subsection{Self-twist rigidity of the branch system}

\begin{theorem}[self-twist rigidity]\label{thm:self-twist}
Let $\mathcal V_\tau$ be the rank-two local system on
$U=\mathbb P^1\setminus\{0,\alpha,\beta,\infty\}$ ($\alpha\beta=1$ the roots
of $q$) defined by the branch operator
$\mathcal D_\tau=4q(t)\theta^2+4t(t-17)\theta+t(t-10)$ annihilating
$\tau=\sqrt F$, with Riemann scheme exponents $\{0,0\}$ at $0$,
$\{0,\tfrac12\}$ at $\alpha,\beta$, $\{\tfrac12,\tfrac12\}$ at $\infty$, and
$\det\mathcal V_\tau=\delta$ the quadratic Kummer system of $\sqrt{q}$-type.
Then, with a fixed determinant normalization, every geometric finite-order
Kummer self-twist of $\mathcal V_\tau$ is trivial, and the conjugate-twist
list is exactly $\{1,\delta\}$; comparing raw and determinant-one
normalizations enlarges the list only to
$\langle\eta\rangle=\{1,\eta,\delta,\eta^{-1}\}$ of order $4$
($\eta=q^{1/4}$).  Moreover $\operatorname{Sym}^2\mathcal V_\tau$ admits no
nontrivial geometric self-twist, the external product on $U\times U$ is
irreducible with only factorwise quadratic twists, and the diagonal rank-four
tensor decomposes only as the Clebsch--Gordan $3+1$.  In particular no Kummer
character of order $>4$ is a self-twist or decomposition character.
\end{theorem}

\begin{proof}[Proof sketch and certification]
Local exponents by direct indicial computation (machine-verified together
with the operator identity $\mathcal D_\tau\tau=0$ to $O(t^{40})$ and the
half-integer recurrences: \texttt{q6372\_local\_system\_check.py}); a
self-twist $L$ must have $L^2$ trivial with ramification supported on the
four singular points, and the exponent lists above eliminate all nontrivial
possibilities via the rigidity lemma (ramification comparison at each cusp);
the tensor statements follow from Zariski density of the connected monodromy
in $SL_2$ and Clebsch--Gordan.  Full proof: working note
\texttt{GPT\_Q6372\_self\_twist\_classification.md} (\S4--\S8), including the
caveat that $\boxtimes$ on $U\times U$ is irreducible and the familiar $3+1$
splitting refers to its diagonal pullback only.
\end{proof}

\noindent This closes the geometric prerequisite of the counting lemma
below: all self-twist and tensor-decomposition exceptions lie in a fixed set
of order $\le4$, so no high-order Mellin character can correlate with the
column event through a hidden twist.

\subsection{The parity law for the zero fibre}

\begin{proposition}[parity law]\label{prop:parity}
Let $Z_p=\{1\le r\le p-2:\ p\mid b_r\}$.  Then
$|Z_p|\equiv\mathbf 1\bigl[p\mid b_{(p-1)/2}\bigr]\pmod 2$.  In particular
$|Z_p|$ is even except at the non-ordinary primes of the level-8 weight-4
newform $\eta(2z)^4\eta(4z)^4$ (Beukers' congruence), of which the only
examples below $30000$ are $p=11,3137$.
\end{proposition}

\begin{proof}
Strong reflection (Theorem~\ref{thm:strong-reflection}) gives the involution
$r\mapsto p-1-r$ on $Z_p$, whose unique fixed point in range is the midpoint
$r=(p-1)/2$.  (Empirical profile across $p<30000$, all $3242$ window primes:
$|Z_p|=2\cdot\mathrm{Poisson}(\lambda\approx0.51)+\text{midpoint indicator}$,
$\chi^2=2.47$ on 3 d.f., mean $1.0185$, max $8$; scan banked in the appendix-Q
ledger.)
\end{proof}

\noindent\emph{Data point.}  The midpoint-zero primes $11$ and $3137$ were
identified by Zudilin (arXiv:2409.00384), who verified that they are the only
examples below $2\times10^4$.  Our scan — two independent C implementations
in exact agreement (78{,}462 window zeros compared bitwise) — extends this:
\emph{$11$ and $3137$ remain the only midpoint-zero primes below $10^6$}
(all $78{,}495$ window primes; a fifty-fold extension of the published
record).  Per-prime zero counts along the way: mean $0.9996$, maximum $12$
(at $p=159977$), in exact agreement with the
$2\cdot\mathrm{Poisson}(0.4998)+\text{midpoint}$ law
($\chi^2=3.43$ on 5 d.f.).

\subsection{The character-order counting lemma and its consequence}

\begin{lemma}[character-order counting]\label{lem:counting}
For $n$ fixed and window primes $p\in(n/2,n]$, the Mellin character of the
column event is the global specialization $\omega_p^{\,n-1}$, of order
$m_p=(p-1)/\gcd(p-1,n-1)$, and
$\#\{p:\ m_p\le T\}\le T\,\tau(n-1)$.
\end{lemma}

\noindent Consequently every bounded-character-order law — including the
apparition theorem (order dividing $8$) and any CM/motivic structure — is
confined to $O(T\tau(n-1))$ window primes and cannot produce a bad column:
$H(n)=H_{\mathrm{high}}(n;T)+O(T\tau(n-1))$.  The pointwise Problem~3.2 is
thereby equivalent to zero density of the \emph{high-order} Mellin bad
diagonal (the canonical frontier statement; see the session notes \S34, \S39).
